%%%%%%%%%%  scan idx 521  |  folio 508  %%%%%%%%%%
\SGAPrintPage{508}%

\sgaStableTarget{sga6.e09.prop.2-2}\textbf{Proposition 2.2.} \uline{With the notation and hypotheses of} 2.1, \uline{let} $f:P\to S$ \uline{be the structural morphism of} $P$. \uline{Then the homomorphism}:
\[
f^*:K_\bullet(S)\longrightarrow K_\bullet(P)
\]
\uline{is surjective}, \uline{and its kernel is the subgroup of} $K_\bullet(S)$ \uline{generated by the subgroups} $(\eta(m)-1)K_\bullet(S)$ \uline{for} $m\in M$, \uline{or}, \uline{equivalently}, \uline{for} $m$ \uline{running through a system of generators of} $M$.

First suppose $r=1$, i.e.\ $T=G_{m\,S}$. Then $M=\mathbb Z$, and for every $m\in M$ one has:
\[
\underline{A}_m\simeq \underline{A}_1^{\otimes m}.
\]
Consequently, putting $\underline{A}_1=\underline{L}$, one has $\underline{A}=\bigoplus_{n\in\mathbb Z}\underline{L}^{\otimes n}$, which shows that $P$ identifies with the complement of the zero section in the bundle $X=\mathbb V(\underline{L})$. Considering $S$ as a closed subscheme of $X$ by the zero section, one has, by 1.1, the exact sequence:
\[
K_\bullet(S)\xrightarrow{g_*}K_\bullet(X)\longrightarrow K_\bullet(P)\longrightarrow 0
\]
where $g:S\to X$ is the zero section of $X$. Moreover, $f^*$ factors as \(K_\bullet(S)\xrightarrow{f'^*}K_\bullet(X)\longrightarrow K_\bullet(P)\), where $f'$ is the structural morphism of $X$. By 1.6, $f'^*$ is an isomorphism, inverse to $g^*$. Thus one has the exact sequence:
\[
K_\bullet(S)\xrightarrow{g^*\circ g_*}K_\bullet(S)\xrightarrow{f^*}K_\bullet(P)\longrightarrow 0,
\]
from which one obtains, by 1.8:
\[
K_\bullet(P)\simeq K_\bullet(S)/\lambda_{-1}(L)K_\bullet(S)=K_\bullet(S)/(1-L)K_\bullet(S)
\]
where $L$ is the class in $K^{\bullet}(S)$ of $\underline{L}$. Since for every $n\in\mathbb Z$, $L^n-1$ is a multiple of $1-L$, this proves 2.2 when $P$ has rank $1$.


%%%%%%%%%%  scan idx 522  |  folio 509  %%%%%%%%%%
\SGAPrintPage{509}%

Continue the proof by induction on the rank $r$ of $P$. Suppose the result proved for torsors under split tori of rank $<r$. Write $M=M_1\oplus M_2$, with $M_1\simeq\mathbb Z^{r-1}$ and $M_2\simeq\mathbb Z$; put $\underline{A}'=\bigoplus_{m\in M_1}\underline{A}_m$ and $\underline{A}''=\bigoplus_{m\in M_2}\underline{A}_m$, so that $\underline{A}=\underline{A}'\otimes_{\underline{O}_S}\underline{A}''$. Let $P_1=\Spec(\underline{A}')$; then $P_1$ is a torsor under $D_S(M_1)$, and $P$ is a torsor over $P_1$ with group $\underline{G}_{m\,P_1}$ and algebra the inverse image of $\underline{A}''$ on $P_1$. By the preceding argument, $K_\bullet(P)$ is the quotient of $K_\bullet(P_1)$ by the subgroup generated by the $(\eta(m_2)-1)K_\bullet(P_1)$ for $m_2\in M_2$. By the induction hypothesis, $K_\bullet(P_1)$ is the quotient of $K_\bullet(S)$ by the subgroup generated by the $(\eta(m_1)-1)K_\bullet(S)$ for $m_1\in M_1$. Taking account of the relation
\[
\eta(m_1+m_2)-1=\eta(m_1)\eta(m_2)-1
=\eta(m_1)(\eta(m_2)-1)+\eta(m_1)-1,
\]
$K_\bullet(P)$ is therefore the quotient of $K_\bullet(S)$ by the subgroup generated by the $(\eta(m)-1)K_\bullet(S)$ for $m\in M$, or equivalently for $m$ running through a system of generators of $M$.

\sgaStableTarget{sga6.e09.cor.2-3}\textbf{Corollary 2.3.} \uline{With the notation and hypotheses of} 2.1, \uline{suppose that} $S$ \uline{is regular}; \uline{then one has a canonical isomorphism}:
\[
K^{\bullet}(P)\simeq K^{\bullet}(S)/\bigl((\eta(m)-1)_{m\in M}\bigr)
\]
\uline{where} $((\eta(m)-1)_{m\in M})$ \uline{is the ideal of} $K^{\bullet}(S)$ \uline{generated by the} $\eta(m)-1$ \uline{for} $m$ \uline{running through a system of generators of} $M$.

If $S$ is regular, the canonical homomorphism $K^{\bullet}(S)\to K_\bullet(S)$ is an isomorphism (IV 2.5). Likewise, $P$ is then regular, so $K^{\bullet}(P)\to K_\bullet(P)$ is an isomorphism. The corollary therefore follows immediately from 2.2.


%%%%%%%%%%  scan idx 523  |  folio 510  %%%%%%%%%%
\SGAPrintPage{510}%

\sgaStableTarget{sga6.e09.sec.3}\subsection*{3. Projective bundles and flag bundles}

\sgaStableTarget{sga6.e09.prop.3-1}\textbf{Proposition 3.1.} \uline{Let} $S$ \uline{be a noetherian scheme}, $\underline{E}$ \uline{a locally free} $\underline{O}_S$\uline{-Module of rank} $n$, $\pi=(p_1,\ldots,p_k)$ \uline{a sequence of integers such that} $\sum_i p_i=n$, \uline{and} $X=\mathbb D_\pi(\underline{E})$ \uline{the scheme of flags of type} $\pi$ \uline{in} $\underline{E}$. \uline{Then the canonical homomorphism}:
\[
K_\bullet(S)\otimes_{K^{\bullet}(S)}K^{\bullet}(X)\longrightarrow K_\bullet(X)
\]
\uline{is an isomorphism}.

Resume the notation of VI 4.5. Thus $\underline{L}_i$ denotes the canonical quotients of the inverse image of $\underline{E}$ on $X$, $\underline{F}_i$ the kernel of $\underline{L}_i\to \underline{L}_{i-1}$ ($i=1,\ldots,k$), and $\lambda_i^{(j)}$ the class in $K^{\bullet}(X)$ of $\wedge^i(\underline{F}_j)$. Then the $\lambda_i^{(j)}$ generate $K^{\bullet}(X)$ as a $K^{\bullet}(S)$-algebra (VI 4.6). Likewise, the images of the $\lambda_i^{(j)}$ in $K^{\bullet}(X_s)$ generate $K^{\bullet}(X_s)$ as an algebra over $K^{\bullet}(k(s))=\mathbb Z$ for every point $s$ of $S$. Since $X_s$ is a flag scheme over a field, it is a regular scheme, so the homomorphism $K^{\bullet}(X_s)\to K_\bullet(X_s)$ is an isomorphism, and thus the homomorphism $K^{\bullet}(X)\to K_\bullet(X_s)$ is surjective. By 1.3, this implies that the homomorphism:
\[
K_\bullet(S)\otimes_{K^{\bullet}(S)}K^{\bullet}(X)\longrightarrow K_\bullet(X)
\]
is surjective.

To prove injectivity, one may first reduce to the case of the complete flag bundle $\mathbb D(E)$. Indeed, suppose that
\[
K_\bullet(S)\otimes_{K^{\bullet}(S)}K^{\bullet}(\mathbb D(\underline{E}))\longrightarrow K_\bullet(\mathbb D(\underline{E}))
\]
is injective. Since $K_\bullet(S)\to K_\bullet(\mathbb D(\underline{E}))$ factors through $K_\bullet(X)$, this

